\documentclass[11pt]{article}
\usepackage{ochamath}
\subjectcolor{2F5DA8}
\setsheet{線形代数・電気系院試補充}{リアプノフ方程式　演習}{https://university-math-crj.pages.dev/linear-algebra/lyapunov-equations/}
\begin{document}
\makesheethead{解答}
自作問題。本文の例題とは数値や設定が異なります。条件・途中式・検算を示してください。
\problem[連続時間の成分計算]
$A=\begin{pmatrix}-2&1\\0&-3\end{pmatrix},Q=I$ のリアプノフ方程式を解き、正定値を確認せよ。
\begin{ans}
$P=\begin{pmatrix}a&b\\b&c\end{pmatrix}$ を入れると $-4a=-1,a-5b=0,2b-6c=-1$。よって $a=1/4,b=1/20,c=11/60$。行列式は $11/240-1/400=13/300>0$、左上も正で正定値。$\dot V=-\|x\|^2$ により漸近安定。
\end{ans}
\point{方程式を解いた後に正定値性を検証する。}
\problem[積分表示]
安定な $A$ に対して $P=\int_0^\infty e^{A^{\mathsf T}t}Qe^{At}dt$ がリアプノフ方程式を満たすことを示せ。
\begin{ans}
安定性から指数関数は減衰し積分は収束。被積分行列を微分すると $A^{\mathsf T}$ の左積と $A$ の右積が現れる。積分して $A^{\mathsf T}P+PA=[e^{A^{\mathsf T}t}Qe^{At}]_0^\infty=-Q$。$Q\succ0$ なら非零 $x$ の $x^{\mathsf T}Px$ は正の積分で正。
\end{ans}
\point{無限遠の境界項が0になる理由を述べる。}
\newpage
\problem[逆向きの安定性]
$P,Q\succ0$ と $A^*P+PA=-Q$ から全固有値の実部が負と証明せよ。
\begin{ans}
$Av=\lambda v$ に左から $v^*$、右から $v$ を使うと $2\operatorname{Re}\lambda\,v^*Pv=-v^*Qv$。非零固有ベクトルでは両二次形式が正なので $\operatorname{Re}\lambda<0$。実行列でも複素固有ベクトルを使えるよう共役転置で書く。
\end{ans}
\point{対角化の仮定は必要ない。}
\problem[離散時間]
$A=\operatorname{diag}(1/3,1/2),Q=I$ に対し $A^{\mathsf T}PA-P=-I$ を解け。
\begin{ans}
対角成分は $((1/3)^2-1)p_{11}=-1$ と $((1/2)^2-1)p_{22}=-1$。非対角成分は $(1/6-1)p_{12}=0$。従って $P=\operatorname{diag}(9/8,4/3)$。正定値で $V(x_{k+1})-V(x_k)=-\|x_k\|^2$。
\end{ans}
\point{連続時間の式と混ぜない。}
\newpage
\problem[Sylvester方程式]
$A=\operatorname{diag}(1,2),B=\operatorname{diag}(3,4),C=\begin{pmatrix}4&10\\15&24\end{pmatrix}$ の $AX+XB=C$ を解け。
\begin{ans}
各成分は $(a_i+b_r)x_{ir}=c_{ir}$。固有値和は4,5,5,6で全て非零なので唯一解。$X=\begin{pmatrix}1&2\\3&4\end{pmatrix}$。左積と右積を足せば元の $C$ が得られる。和が0になる成分があれば右辺によって解なしまたは非一意となる。
\end{ans}
\point{一意性を固有値の和で先に確認する。}
\problem[解はあっても正定値でない]
スカラー $A=2,Q=1$ で連続時間のリアプノフ方程式を解き、安定性との関係を述べよ。
\begin{ans}
$2P+2P=-1$ から $P=-1/4$。方程式は唯一解を持つが正定値ではない。系 $\dot x=2x$ は不安定で、正定値解の存在という安定性の定理の条件を満たさない。固有値和が非零という条件は代数的一意性だけを保証する。
\end{ans}
\point{解の存在と正定値解の存在を区別する。}
\end{document}
